% =========================================================================
% English translation (generated by AI using Claude Opus 4.8 (Max effort)).
% This file was translated from Hungarian to English by AI.
% DISCLAIMER: Please review against the original Hungarian .tex files, before relying on it!
% SOURCE: 16. Haladó algoritmusok.tex
% =========================================================================
\documentclass[margin=0px]{article}

\usepackage{listings}
\usepackage[utf8]{inputenc}
\usepackage{graphicx}
\usepackage{float}
\usepackage[a4paper, margin=0.7in]{geometry}
\usepackage{subcaption}
\usepackage{amsthm}
\usepackage{amssymb}
\usepackage{amsmath}
\usepackage{fancyhdr}
\usepackage{setspace}

\onehalfspacing

\renewcommand{\figurename}{Figure}
\newenvironment{tetel}[1]{\paragraph{#1 \\}}{}
\newcommand{\R}{\mathbb{R}}

\pagestyle{fancy}
\lhead{\it{CS BSc Final Exam Topics}}
\rhead{16. Advanced algorithms}

\title{\textbf{{\Large ELTE IK - Computer Science BSc} \vspace{0.2cm} \\ {\huge Final Exam Topics}} \vspace{0.3cm} \\ 16. Advanced algorithms}
\author{}
\date{}

\begin{document}
\maketitle

\begin{center}
\fbox{\parbox{0.92\textwidth}{\small \textit{Note: This document was translated from Hungarian to English by AI. Please verify the information against the original Hungarian source before relying on it.}}}
\end{center}

\begin{tetel}{Advanced algorithms}
    Elementary graph algorithms: breadth-first, depth-first traversal and its applications. Minimum spanning trees, a general algorithm, Kruskal's and Prim's algorithms. Single-source shortest paths, queue-based Bellman-Ford, Dijkstra, DAG shortest path. All-pairs shortest paths: Floyd-Warshall algorithm. Transitive closure of a graph.
\end{tetel}

\section{Graph algorithms}
\subsection{Graph representation}
\begin{description}
    \item[Linked-list representation] \hfill \\
        Let us place the vertices of the graph in an array (or linked list). To every element let us assign a linked list in which we list the neighbors of the given vertex (with the possible edge weights).
    \item[Matrix representation] \hfill \\
        Let the number of vertices be $n$. Then in a matrix $A^{n\times n}$ we mark which vertices are connected. Then both in the rows and in the columns the vertices appear, and in cell $a_{ij}$ the weight of the edge leading from vertex $i$ to vertex $j$ appears; if there is no edge between the two vertices, then $-\infty$ (in the unweighted case $1$ and $0$)

        If the graph is undirected, obviously $a_{ij} = a_{ji}$
\end{description}
\subsection{Breadth-first traversal}
We reach the vertices in the order of distance from the start vertex $s$ of a (directed/undirected) graph $G$. It gives the spanning tree of the shortest paths, so only the distance matters, not the weight.

We store the tracked vertices in a queue data structure; we put the children of the current vertex into the queue. The next vertex will be the very first element of the queue.

We write the distances of the vertices into a $d$ array, their parents into a $\pi$ array, and initialize with $\infty$ and $0$ respectively.
\begin{flushleft}
The algorithm:
\end{flushleft}
\begin{enumerate}
    \item We put the start vertex $s$ into the queue
    \item We repeat the following steps until the queue is empty
    \item We take out the very first ($u$) element of the queue
    \item Those child vertices whose distance is not $\infty$ we disregard (we have already been on these)
    \item For the other children ($v$): we set their parent ($\pi[v] = u$), and their distance ($d[v] = d[u]+1$). Then we put them into the queue.
\end{enumerate}
Its applications
\begin{enumerate}
	\item Finding the shortest path between two nodes $u$ and $v$, measuring the length of the path by the number of edges (advantage compared to depth-first search)
	\item (Reverse) Cuthill–McKee mesh numbering, a variant of breadth-first traversal
	\item Serialization/deserialization of a binary tree versus serialization in sorted order; it enables the efficient reconstruction of the tree
	\item Testing the bipartiteness of a bipartite graph
\end{enumerate}
\subsection{Searching for minimum-cost paths}
\begin{description}
    \item[Dijkstra algorithm] \hfill \\
        In a directed or undirected graph $G$ with positive edge weights, it searches for minimum-cost paths from the start vertex $s$ to every vertex.

        The algorithm is a modified version of breadth-first traversal. Since here a longer path can have a smaller cost than a shorter one, an already-found vertex must not be disregarded. Therefore every vertex has three states (not reached, reached, done). We handle the $d$ and $\pi$ arrays similarly to breadth-first traversal.

        The not-yet-done vertices we place into a priority queue, or in every case we apply minimum search.

        The algorithm:
        \begin{enumerate}
            \item We set the weight of the start vertex $s$ to 0 and store it
            \item We repeat the following steps until our container is empty
            \item We take out the best ($u$) element of the queue, and set it to ``done''
            \item If a child vertex ($v$) is not done, and the weight of the path currently leading to it is smaller than the previous one, then: we set its parent to $u$ ($\pi[v] = u$), and update its weight ($d[v] = d[u]+d(u,v)$).
            \item The other vertices we skip.
        \end{enumerate}

    \item[Bellman-Ford algorithm] \hfill \\
        From the start vertex $s$ of an edge-weighted (even negative) directed graph $G$ it searches for minimum-cost paths to every edge, and recognizes if there is a negative-cost cycle in the graph. We handle the $d$ and $\pi$ arrays similarly to the previous ones.

        The algorithm:
        \begin{enumerate}
            \item Let us set the weight of the start vertex to 0.
            \item In $n-1$ iterations let us go through all vertices, and compare every vertex ($u$) with every vertex ($v$). If we found a cheaper path then in $v$ we overwrite its weight ($d[v] = d[u]+d(u,v)$), and its parent ($\pi[v] = u$).
            \item If in the $n$-th iteration there was also a modification, there is a negative cycle in the graph
        \end{enumerate}
\end{description}
\subsection{Searching for a minimum-cost spanning tree}
The minimum-cost spanning tree or minimum spanning tree is the spanning tree of smallest edge weight in a connected, undirected graph. A spanning tree is a tree that contains all vertices of the graph and whose edges are among the edges of the original graph. The minimum spanning tree is not necessarily unique, but its weight is. For searching for any minimum spanning tree of a graph, Kruskal's and Prim's algorithms can be used. Both algorithms use a greedy strategy.
\begin{description}
	\item[Kruskal algorithm] \hfill \\
	   We sort the edges in increasing order of their weight, and examine in turn which ones we may include in the solution. Initially every vertex of the graph forms a separate set. An examined edge gets included in the solution if its two endpoints are in two different sets, because from this we know that adding the examined edge does not create a cycle in the graph. Then we merge the two sets. At the end of the algorithm a single set will remain.

	  The algorithm:
	  \begin{enumerate}
		 \item Let us choose the edge of smallest weight.
		 \item If adding the edge to the subgraph forms a cycle, let us discard it, otherwise let us add it.
		 \item Let us repeat the above steps until there is an unexamined edge.
	  \end{enumerate}
	  Its applications:
	  \begin{enumerate}
	  	 \item The Kruskal algorithm was made for the simple proof of the theorem that the minimum-cost spanning tree is unique if there are no two edges of equal weight in the graph.
	  	 \item A random maze can be generated with it. In this case the weight of every edge of the graph to be processed coincides, so its edges do not have to be sorted.
	  \end{enumerate}
	\item[Prim algorithm] \hfill \\
	The Prim algorithm searches for a minimum-cost spanning tree from the start vertex $s$ of an undirected edge-weighted (even negative) graph. We handle the $d$ and $\pi$ arrays similarly to the previous ones. The algorithm places the vertices into a priority queue.

	The algorithm:
	\begin{enumerate}
	    \item Let us set the weight of the start vertex to 0.
	    \item We place the vertices into the priority queue.
	    \item We perform the following steps until the priority queue is emptied.
	    \item We take out a vertex ($u$) from the queue.
	    \item For every child of it ($v$) that is still in the queue and the weight of the recorded edge leading to $v$ is greater than the one just found: We change the parent of $v$ to $u$, overwrite the recorded weight $d[v] = d(u,v)$. Then we overwrite the state of $v$ in the priority queue.
	    \item With those children that are not in the queue, or whose weight we cannot improve, we do not change anything.
	\end{enumerate}

$\textbf{ALTERNATIVELY}$

Its operating principle is that it builds the tree vertex by vertex, starting from an arbitrary vertex and at every step finding the cheapest possible edge to a next vertex.

	  The algorithm:
	  \begin{enumerate}
		\item Let us choose an arbitrary vertex of the graph; let this be a single-vertex tree.
		\item As long as there is a vertex of the original graph that is not in the tree, let us repeat the following steps.
		\item Let us choose the one of smallest weight among the edges running between the vertices of the tree and the other vertices of the graph.
		\item Let us move the non-tree vertex of the chosen edge into the tree together with the edge.
	  \end{enumerate}
	  Application: This too can be used for generating a random maze.
\end{description}
\subsection{Depth-first traversal}
By the depth-first traversal of a directed (not necessarily connected) graph $G$ we obtain a depth-first tree (forest). The algorithm is the following:
\begin{itemize}
    \item The edge weights play no role
    \item There is no start vertex; we start the algorithm for every vertex of the graph. (Of course, then, if we choose a vertex we have already been on, the algorithm does not start.)
    \item We choose the vertices greedily, that is, choosing the first among the children of every vertex, we proceed forward as long as possible. (We find a vertex that has no child, or whose every child we have already been on.)
    \item If we can no longer proceed forward, we step back.
    \item To every vertex we assign two values. One is the depth-first number, which marks how many-th we reached it. The other is the finishing number, which marks how many-th we stepped back from it.
\end{itemize}
The edges of the graph can be classified during the depth-first traversal. (At initialization we set every value to 0)
\begin{itemize}
    \item Tree edge: The depth-first number of the next vertex is 0
    \item Back edge: The depth-first number of the next vertex is greater than 0, and its finishing number is 0 (So we turn back to a previous vertex of the current path)
    \item Cross edge: The depth-first number of the next vertex is greater than 0, and its finishing number is also greater than 0, and furthermore the depth-first number of the current vertex is greater than the depth-first number of the next vertex. (Then we are dealing with an edge starting from a vertex preceding the current vertex and pointing into an already-found path)
    \item Forward edge: The depth-first number of the next vertex is greater than 0, and its finishing number is also greater than 0, and furthermore the depth-first number of the current vertex is smaller than the depth-first number of the next vertex. (Then we are dealing with an edge starting from the current vertex and pointing into an already-found path)
\end{itemize}
Algorithms that use depth-first traversal as a building block
\begin{enumerate}
	\item Finding the connected component of a graph
	\item Topological sorting
	\item Finding 2- (edge or vertex) connected components
	\item Finding 3- (edge or vertex) connected components
	\item Finding separating edges in a graph
	\item Finding the strongly connected components of a graph
	\item Solving puzzles with only one solution, for example mazes
	\item The creation of a maze can use randomly performed depth-first search
\end{enumerate}
\subsection{Topological sorting of a DAG}
\begin{description}
    \item[Basic concepts] \hfill
        \begin{itemize}
            \item Topological sorting: \\
                  A topological sorting of a graph $G(V,E)$ is an order of the vertices in which for $\forall (u\rightarrow v) \in E$ edge $u$ is earlier in the order than $v$
            \item DAG - Directed Acyclic Graph: \\
                  A directed acyclic graph. \\
                  Most often used for directing workflows and analyzing dependencies.

                  Properties:
                  \begin{itemize}
                      \item If for a graph $G$ the depth-first traversal finds a back edge (that is, found a cycle) $\Longrightarrow$ $G$ is not a DAG
                      \item If $G$ is not a DAG (there is a cycle in it) $\Longrightarrow$ Any depth-first traversal finds a back edge
                      \item If $G$ has a topological sorting $\Longrightarrow$ $G$ is a DAG
                      \item Every DAG can be topologically sorted.
                  \end{itemize}
        \end{itemize}
    \item[Topological sorting of a DAG] \hfill \\
        During the depth-first traversal of a graph $G$ let us put into a stack those vertices from which we stepped back. After the algorithm, writing out the content of the stack, we obtain a topological sorting of the graph.
\end{description}
\subsection{All-pairs shortest paths: Floyd-Warshall algorithm}

We write the distances of the vertices into a $d$ matrix, their parents into a $\pi$ matrix, which we initialize, similarly to the previous ones, with $\infty$ and 0 values. The end result will be that $d[i, j]$ is the length of the optimal path leading from the vertex with index i to the vertex with index j, or $\infty$ if there is no path from i to j. And $\pi[i, j]$ is, computed by the algorithm, the direct predecessor of vertex j on the optimal path leading from the vertex with index i to the vertex with index j, if i $\ne$ j and a path from i to j exists, otherwise $\pi[i, j]$ = 0.
$\textbf{Description of the algorithm:}$

Take a graph $G$, with vertices $V$ numbered from 1 to $N$. Furthermore a function, the so-called $\textbf{shortestPath(i,j,k)}$, which returns the shortest path between $i$ and $j$ using given vertices: ${1,2,...,k}$ as intermediate points along the vertices. Taking the given function into account, our goal is to find the shortest path from every $i$ to $j$ using any of the vertices ${1,2,...,N}$.
For each of the vertex pairs, the $\textbf{shortestPath(i,j,k)}$ can be either\\
(1) a path that does not pass through $k$ (which uses only the given vertices: ${1,...,k-1}$)\\
OR\\
(2) a path that passes through $k$ (from $i$ to $k$ and then from $k$ to $j$, in both cases using the given intermediate vertices: ${1,...,k-1}$).

We know that the best path from $i$ to $j$ that uses only those vertices which are from $1$ through $k-1$ is determined by the $\textbf{shortestPath(i,j,k-1)}$, and it is clear that if there were a better path from $i$ to $k$ and then from there to $j$, then the length of these paths would be the concatenation of the shortest path from $i$ to $k$ (using only the intermediate vertices ${1,...,k-1}$) and the shortest from $k$ to $j$ (using only the intermediate vertices ${1,...,k-1}$).

If $\omega(i,j)$ is the weight of the edge between the given vertices $i$ and $j$, then we can determine the $\textbf{shortestPath(i,j,k)}$ function according to the following recursive formula:
$\textbf{shortestPath(i,j,0)}$ = $\omega(i,j)$
recursive case:\\
$\textbf{shortestPath(i,j,k) =}$\\
$\textbf{min(shortestPath(i,j,k-1), shortestPath(i,k,k-1)}$ + $\textbf{shortestPath(k,j,k-1)}$).

This formula is the heart of the Floyd-Warshall algorithm. During the run of the algorithm it first computes, based on\\
$\textbf{shortestPath(i,j,k)}$, the $(i,j)$ pairs for the cases $k = 1$, then $k = 2$ and so on. This process continues until finally $k = N$ holds and the N-th iteration also runs, and we have found the shortest path for each $(i,j)$ pair using any intermediate vertex.

\subsection{Transitive closure of a graph}
$\textbf{First definition:}$ By the transitive closure of the graph $G = (V, E)$ we mean the relation V × V $\supseteq$ T, where (u, v) $\in$ T if in $G$ the vertex v is reachable from the vertex u.

\begin{flushleft}
$\textbf{Second definition:}$ The transitive closure of a directed graph $G = (V,E)$ is the graph $G' =(V,E')$ in which an edge leads from u to v exactly when a directed path leads from u to v in the original graph $G$.\\
\end{flushleft}

To every directed acyclic graph ($\textbf{DAG}$) a partial order of its vertices can be assigned, in which u $\le$ v holds exactly when in the graph there exists a directed path from u to v. At the same time, such a partial order can be represented by many different directed acyclic graphs. Of these, the transitive reduction contains the fewest edges, and the transitive closure the most.

The transitive closure of a graph can easily be computed, e.g., with the Floyd-Warshall algorithm in O($n^{3}$) time, or with n breadth-first searches in O(n(n+m)) time. At the same time there are also significantly more efficient methods that in practice run in almost linear time.

One such method is based on the determination of the strongly connected components, as well as on the topological sorting of the graph of the components, but it also applies complex data structures.
\end{document}
